\chapter*{Preface}
\addcontentsline{toc}{chapter}{Preface}

This book is about one question and one family of examples. The question is
the Hodge conjecture: on a smooth complex projective variety, is every
rational cohomology class of type $(p,p)$ a rational combination of classes
of subvarieties? The examples are abelian varieties, and among them the
abelian varieties of Weil type. That is where the first Hodge classes appear
that divisors cannot produce.

The conjecture is not proved here. Nobody has proved it, and a reader who
opens the book hoping for a proof should know that at the start. What the
book does is narrower and, I hope, more useful. It explains why the
conjecture is true where it is known to be true, with complete proofs where
the arguments are short enough to give. It explains what changes when one
moves from divisors to classes of higher codimension. And it identifies, as
precisely as current knowledge allows, the statement that is still missing
for abelian varieties. For a general abelian variety of Weil type the whole
conjecture comes down to the algebraicity of a single explicit class. That
reduction is classical in spirit, since it goes back to Weil, but its
consequences are worth following carefully, and most of the book does so.

The first three chapters are background. Chapter~\ref{ch:conjecture} states
the conjecture, explains why the hypotheses are what they are, and lists the
counterexamples to its stronger variants. Chapter~\ref{ch:toolkit} collects
the Hodge theory the rest of the book uses. Chapter~\ref{ch:lefschetz}
proves the theorem of Lefschetz on $(1,1)$ classes, which settles degree
two. The proof is short, and the way it fails in higher degree explains a
good deal of the difficulty.

Chapters~\ref{ch:abelian} and~\ref{ch:weil} are the core of the book. The
cohomology of an abelian variety is an exterior algebra, its Hodge classes
are the invariants of a reductive group, and the question becomes one of
representation theory followed by one of geometry. Weil's classes are the
first case where the representation theory says yes and the geometry has
not caught up. Chapter~\ref{ch:known} records the cases in which it has:
the work of Schoen, Koike and Markman, the members where divisors
suffice, and the more recent preprints, kept clearly separate from refereed
theorems. Chapter~\ref{ch:remaining} is an attempt at the
remaining case. It follows the natural route, deforming a cycle from a
special member of the family to a general one, and shows exactly where that
route stops.

The reader I have in mind knows the basic language of algebraic geometry
and complex manifolds at the level of Griffiths and Harris or Voisin's two
volumes. Where a proof would need more than a few pages of machinery that
is well covered elsewhere, I state the result precisely, give a reference,
and say what the proof uses.

Some of the material on Weil classes and their algebraic loci first
appeared in my preprint \emph{Algebraic Loci of Weil Classes on Abelian
Varieties} \cite{BhH8}. Most of what is taken from it is proved again here.
Where a proof is only sketched, the text says so and points to the
preprint. The arguments are analytical
throughout: no step rests on a computer calculation. Every place where a
result depends on an unrefereed preprint by someone else is marked as such.

\subsection*{A note on preparation}
Drafts of parts of this book were prepared with the help of an AI language
model, used to draft and revise text under the author's direction. Every
statement and proof has been checked by the author, who is responsible for
the content.

\bigskip
\noindent Deep Bhattacharjee
